\section{Acyclic resolutions with a cardinal bound}
\label{sec:A}

All groups, simplicial complexes and CW complexes in this section belong to
$\mathcal U$. Mapping spaces are derived, unpointed mapping spaces. A map
is called acyclic when every homotopy fibre is nonempty, connected and has
zero reduced integral homology. Equivalently it induces isomorphisms in
homology with every local system on its target. Base points used in writing
fundamental groups do not enter the unpointed conclusions.

\begin{claim}{A.1}{Proved}{The simple acyclic envelopes used in R2}
For every infinite cardinal $\lambda\in\mathcal U$ and group $H$ with
$|H|\leq\lambda$, there is an embedding $H\hookrightarrow C$ with $C$
nonabelian simple, integrally acyclic, and $|C|=\lambda$.
\end{claim}
\begin{proof}
Here is a check of GALOIS B.10, including the two constructions on which
it rests. Start with $H*F_\lambda$. For a current group $L$ of cardinal
$\lambda$ and every pair $g\ne1,h\in L$, adjoin an independent free letter
$t=t_{g,h}$ and put
\[
 w=g t g^{-1}t^{-1}g t^2g^{-1}t^{-2}.
\]
The reduced free-product word $w$ is cyclically reduced of length eight.
Cyclic reduction of $wh$ gives length eight when $hg\ne1$, and length
six when $hg=1$; in the latter case the two adjacent powers of $t$
combine to a nonzero power. Thus both words have infinite order. Adjoin
an HNN letter $s_{g,h}$ with $s^{-1}ws=wh$. The embeddings of the
associated infinite cyclic groups and the HNN normal form show that
$L$ embeds. The equation $h=w^{-1}s^{-1}ws$ puts $h$ in the normal
closure of $g$. Repeat countably. Every two elements of the union occur
at a common stage, so the resulting group is simple. It contains
$F_\lambda$, is nonabelian, and still has cardinal $\lambda$.

For the other envelope, form $L\times L$, with $j(l)=(l,1)$ and
$k(l)=(1,l)$, and adjoin two HNN letters satisfying
\[
 d^{-1}j(l)d=j(l)k(l),\qquad s^{-1}j(l)s=k(l).
\]
The axes and diagonal are embedded copies of $L$, so again the original
group embeds. Denote this group by $m(L)$. For a field $k$, suppose
$f:H_0\to L$ vanishes on positive homology below degree $n$. The
K\"unneth formula shows that the image of the degree-$n$ diagonal after
$f$ has only its two axis terms. In $m(L)$ the two axes and the diagonal
are conjugate to $j$. Consequently $(jf)_*=2(jf)_*$ in degree $n$,
and $(jf)_*=0$. Induction proves that $n$ successive mitoses kill
homology through degree $n$. In the countable union every bar cycle
has finite support, so positive homology vanishes over every field.
Rational coefficients make every positive integral homology group
torsion, and the universal coefficient sequence with $\mathbb F_p$ in
the next degree excludes its $p$-torsion, for each prime $p$. Integral
homology therefore vanishes as well. The cardinal bound is preserved.

Alternate the simple and acyclic envelopes and take their countable
union. Its cofinal simple stages make it simple; its cofinal acyclic
stages make it acyclic. This confirms the construction used by both
written proofs of R2. The mitosis construction also occurs in
\cite[Sections 4--5]{A_BDH1980}; the normal-form and homology calculations
above supply the particular assertions used here.
\end{proof}

\begin{claim}{A.2}{Proved}{The complete mapping-space assertion A4}
Let $t:X\to W$ be an acyclic map of connected CW complexes, let
$P=\ker\pi_1(t)$, and let $Y$ be any $\mathcal U$-small space, possibly
disconnected or empty. Restriction is a weak equivalence
\[
 t^*:\operatorname{Map}(W,Y)\ \longrightarrow\
 \operatorname{Map}(X,Y)_{P=1},
\]
where the right side is the union of those components represented by maps
$f:X\to Y$ whose induced fundamental-group homomorphism kills $P$.
In particular this holds for every resolution $BG\to W$ in R2.
\end{claim}
\begin{proof}
We check the specified map, including its higher homotopy groups. Given
$f:X\to Y$, first work in the component $Y_0$ containing its image and
form the homotopy pushout $D=Y_0\cup_X^h W$. Write
$j:Y_0\to D$ and $k:W\to D$. Since $\pi_1(t)$ is onto,
van Kampen identifies $\pi_1D$ with the quotient of $\pi_1Y_0$ by the
normal closure of $f_*(P)$. For every local system $M$ on $D$, the
relative Mayer--Vietoris comparison identifies
$H_*(D,Y_0;M)$ with $H_*(W,X;k^*M)$, which vanishes by acyclicity of
$t$. If $f_*(P)=1$, then $j$ is an isomorphism on fundamental groups
and on homology with all local coefficients. Apply the latter assertion
to $\mathbb Z[\pi_1D]$; the resulting map of simply connected universal
covers is an integral homology equivalence, hence a weak equivalence by
Hurewicz and Whitehead. Thus $j$ is a weak equivalence. A homotopy inverse
$r$ gives $rk:W\to Y_0$ and a homotopy $(rk)t\simeq f$.
The reverse implication follows by applying $\pi_1$ to such a factorization.
This proves the assertion about the image on components without assuming
an extension theorem for the plus construction. It is also the pushout
argument in \cite[Proposition 3.1]{A_HH1979}.

Take $Y_0=W$ and $f=t$ in the same calculation. The homotopy codiagonal
$c:W\cup_X^h W\to W$ is a weak equivalence: its fundamental-group map
is the fold of $\pi_1W*_{\pi_1X}\pi_1W$, an isomorphism, and its
local-coefficient Mayer--Vietoris comparison is an isomorphism. Applying
$\operatorname{Map}(-,Y)$ identifies $c^*$ with the diagonal
\[
 \operatorname{Map}(W,Y)\longrightarrow
 \operatorname{Map}(W,Y)
 \mathop{\times}^{h}_{\operatorname{Map}(X,Y)}
 \operatorname{Map}(W,Y).
\]
Therefore every homotopy fibre $F$ of $t^*$ has equivalence diagonal
$F\to F\times F$. A nonempty such $F$ is contractible: its set of
components has bijective diagonal, so is a singleton, and its homotopy
groups have diagonal isomorphisms into their squares, which forces
all of them to vanish. This establishes the weak equivalence on exactly
the previously identified components. The argument applies in each
component of $Y$; when $Y$ is empty both mapping spaces are empty.
\end{proof}

\begin{claim}{A.3}{Proved}{Step-by-step confirmation of R2, Proof 2}
Let $W$ be a connected $\mathcal U$-small CW complex and
$\mu\in\mathcal U$ a cardinal. There is an acyclic map $t:BG\to W$
whose fundamental-group map is onto with perfect kernel $P$, and every
homomorphism $G\to L$ with $|L|\leq\mu$ kills $P$.
If $W$ has a simplicial model with $\nu$ simplices, the construction
uses at most $\max(\aleph_0,\mu^+,\nu)$ cells and a group of at most
that cardinality.
\end{claim}
\begin{proof}
We follow GALOIS B.11, retaining its maps and checking each induction.
Replace $W$ by a homotopy equivalent simplicial complex $K$. Put
$\lambda=\max(\aleph_0,\mu^+)$. In dimensions zero and one set
$T(L)=|L|$ for every subcomplex $L\subseteq K$. Graph components are
aspherical, and a subgraph injects on fundamental groups by reduced
edge-path normal form.

Suppose the construction is made through dimension $n-1$. For each
$n$-simplex $\sigma$, $T(\partial\sigma)$ is connected and aspherical.
Its group has cardinal at most $\lambda$: the boundary has finitely
many faces and each earlier cone has at most $\lambda$ cells. Use A.1
to embed this group in a simple acyclic group $C_\sigma$ of cardinal
$\lambda$. Realize the embedding by a cellular map
$a_\sigma:T(\partial\sigma)\to BC_\sigma$, extending equivariantly
on contractible universal covers. Let $Q_\sigma$ be its mapping
cylinder, with $T(\partial\sigma)$ as free end, and set
\[
 T(L)=T(L^{(n-1)})\cup_{\coprod_{\sigma\in L_n}T(\partial\sigma)}
                       \coprod_{\sigma\in L_n}Q_\sigma.
\]
All choices for a simplex are made once. The map $t_L:T(L)\to |L|$
contracts the image of each boundary linearly toward the barycentre
on its cylinder and sends $BC_\sigma$ to the barycentre. It is an
actual continuous map and commutes with subcomplex inclusions.

For completeness, the asphericity argument concerns a graph of spaces,
with aspherical vertex and edge spaces and injective edge groups.
Amalgam/HNN normal form embeds the vertex groups. In the universal
cover the incidence graph is a tree; the lifted vertex and edge spaces
are contractible. Each finite subtree is contractible by successively
removing terminal vertices. A sphere or a homotopy has compact image
in a finite CW subcomplex and hence projects into a finite subtree.
The universal cover is weakly contractible and therefore contractible.
The ordinary gluing above is equivalent to this mapping-cylinder model
because its attaching inclusions are CW cofibrations.

This argument first gives asphericity of $T(L)$ and injection of its
lower skeleton. To check the stronger subcomplex assertion, for
$L\subseteq K$ of dimension at most $n$ form precisely
\[
 V=T(K^{(n-1)})\cup_{T(L^{(n-1)})}T(L).
\]
For each component of the intersection, both group homomorphisms are
injective, by the preceding assertion and the dimension-$n-1$
induction. The graph-of-spaces calculation embeds both sides in $V$.
Attach next the cones for $K_n\setminus L_n$. Their attaching groups
already inject into $T(K^{(n-1)})$, hence into $V$, and remain injective
at every subsequent attachment. The result is exactly $T(K)$, so
$T(L)\to T(K)$ is injective on fundamental groups. This verifies the
specific intermediate space $V$ used in the written proof. Union over
dimensions preserves all assertions by finite support of spheres,
loops and homotopies.

Let $M$ be any local coefficient system on $|L|$. Over an attached
simplex its pullback is constant. The long exact sequences of the pairs
$(Q_\sigma,T(\partial\sigma))$ and $(\sigma,\partial\sigma)$,
acyclicity of $Q_\sigma$, and the previously proved boundary comparison
show that the specified map is an isomorphism on relative homology
with every coefficient group. Excision and the skeletal long exact
sequence give
\[
 H_*(T(L);t_L^*M)\ \xrightarrow{\cong}\ H_*(|L|;M).
\]
Cellular chains with local coefficients commute with the skeletal
union, so this includes infinite-dimensional $L$.

Let $N$ be the normal closure in $G=\pi_1T(K)$ of all the $C_\sigma$
for $\dim\sigma\geq2$. Van Kampen identifies the specified quotient
$G/N\to\pi_1|K|$. In dimension two it imposes exactly each triangular
boundary relator. In dimension $n\geq3$, the induction applied to
$\partial\sigma$ identifies its group modulo its own lower cone groups
with $\pi_1\partial\sigma=1$; those cone groups are already killed.
Thus there is no additional relation in the quotient. Taking unions
proves $G/N\cong\pi_1|K|$, and $N=\ker\pi_1(t_K)$. Each conjugate
of each $C_\sigma$ is perfect inside $N$, and together they generate
$N$, so $N=[N,N]$. Every homomorphism to a group of cardinal at most
$\mu$ kills every $C_\sigma$, by simplicity and its cardinality, and
therefore kills $N$.

The bar models, cylinders and unions have at most
$\max(\lambda,\nu)$ cells; this also bounds their fundamental groups.
Choosing $BG\simeq T(K)$ and composing with $|K|\simeq W$ proves the
claim. This confirms each part of Proof 2, and improves its bound by
observing that only the standard finite boundaries determine the sizes
of its individual cones. A.2 gives the stated mapping-set consequence.
\end{proof}

\begin{claim}{A.4}{Proved}{The size lemma and the repair of R2, Proof 1}
For every connected CW complex $E$ with at most an infinite cardinal
$\lambda$ of cells, there is an acyclic map $BH\to E$ with
$|H|\leq\lambda$ and a classifying model with at most $\lambda$ cells.
With this lemma, every step of \texttt{OWN\_RESULTS.md}, Theorem 3.1,
is valid. Its recalled bound is supplied here by the later metric
Kan--Thurston theorem, rather than attributed to the original
Kan--Thurston construction.
\end{claim}
\begin{proof}
A CW complex with at most $\lambda$ cells has a homotopy equivalent
simplicial complex with at most $\lambda$ simplices, of the same
dimension. The skeletal construction in
\cite[Theorem 2C.5]{A_Hatcher2002} supplies the precise models: replace
each attaching sphere by a finite triangulation, approximate its map
simplicially, and use the simplicial mapping cylinder and cone defined
there. The target skeleton is retained as a subcomplex. Each cell adds
only finitely many simplices, because its attaching sphere has compact
image in a finite subcomplex. Thus the stated construction gives the
claimed bound for every infinite $\lambda$, beyond the countable case
mentioned explicitly in the theorem.

Apply \cite[Theorem A]{A_Leary2013}. Its stated finite-complex and union
properties mean that the resulting cubical complex is the union of at
most $\lambda$ finite complexes, one for each finite subcomplex of the
input. It is therefore a $\lambda$-cell model. Its components are
aspherical, and its map is an isomorphism on homology with every local
system, exactly the assertion needed. The theorem's connectedness
and surjectivity on fundamental groups are also explicit in its
construction and Remark 6.6. This proves the size lemma.

Here is the audit of the transfinite proof. Well-order the cells of
dimension at least two in nondecreasing dimension; an ordinal of
cardinality at most $\lambda$ suffices. Start with the free group of
the one-skeleton. At a cell attached by $a:S^m\to W_\alpha$, for
$m=1$ use $H'=\mathbb Z$ and a lift of its loop class under the
surjection $G_\alpha\to\pi_1W_\alpha$. For $m\geq2$ form the
homotopy pullback
\[
 E=S^m\mathop{\times}^{h}_{W_\alpha}BG_\alpha.
\]
It is connected and $E\to S^m$ is acyclic, since its fibre is the
acyclic fibre of $BG_\alpha\to W_\alpha$. Its homotopy groups have
cardinality at most $\lambda$ by the two homotopy-fibre exact sequences.
It consequently has a $\lambda$-cell CW model: choose generators and
relations for its fundamental group and then add at most $\lambda$
cells at each stage to realize and kill the higher homotopy groups.
Apply the preceding size lemma to this model. The resulting
$q:BH'\to S^m$ is acyclic, and the other projection is represented by
a homomorphism $u:H'\to G_\alpha$, because its target is aspherical.
The pullback supplies the required homotopy
$t_\alpha Bu\simeq aq$.

Choose simple acyclic groups $A',A''$ of cardinal $\lambda$ containing
$H'$ and form the specified amalgam
\[
 G_{\alpha+1}=(G_\alpha\times A')*_{H'}A'',
 \qquad h\longmapsto (u(h),i'(h)).
\]
Both amalgamating maps are injective; normal form embeds $G_\alpha$.
Use a double mapping cylinder as the classifying model, containing the
previous model as a CW subcomplex. On its first vertex use
$t_\alpha\operatorname{pr}_1$, on its other vertex the centre of the
new cell, and on the cylinders use the pullback homotopy followed by
the cone on $q$. Thus the maps really agree with the previous map,
not merely after passage to homotopy classes.

For a local system on $W_{\alpha+1}$, the $A'$-action on its pullback
is trivial. The projection $BG_\alpha\times BA'\to BG_\alpha$
is a homology equivalence with these coefficients, by the product
spectral sequence and acyclicity of $A'$. The edge comparison is $q$;
the other vertex has constant coefficients and is acyclic. The
Mayer--Vietoris comparison therefore proves acyclicity at the successor
stage. For a homomorphism to a group of size $<\lambda$, both $A'$
and $A''$ have zero image. Their amalgamation forces $u(H')$ to have
zero image too. By the previous factorization property the map factors
through $\pi_1W_\alpha$, and the only new fundamental-group relation
is the attaching loop when $m=1$. For $m\geq2$, $\pi_1S^m=1$, so
$t_{\alpha*}u(H')=1$. In both cases the map factors through
$\pi_1W_{\alpha+1}$.

At a limit ordinal take the union of these actual CW models and group
embeddings. Finite support gives the colimit calculation for homology
with every pulled-back local system. The factorizations through
$\pi_1W_\alpha$ are unique because these quotient maps are onto;
therefore they are compatible and yield the factorization at the
limit. Cardinalities remain at most $\lambda$. These arguments check
the lifting, amalgam, Mayer--Vietoris, quotient and limit steps of
Proof 1. The needed bound is now established, so it remains no
unverified premise in that proof.
\end{proof}

\begin{claim}{A.5}{Proved}{A relative form of the checked construction}
Let $f:BH\to W$ be a based map of connected $\mathcal U$-small CW
complexes, with $BH$ aspherical and $f_*:H\to\pi_1W$ onto. There are
an embedding $j:H\hookrightarrow G$ and an acyclic map $t:BG\to W$
such that $tBj=f$ in a CW model containing $BH$ as a subcomplex.
For an infinite $\lambda$ bounding the cells of $BH$ and $W$, all
models and groups can be chosen of size at most $\lambda$.
\end{claim}
\begin{proof}
Construct a relative CW approximation $BH\subseteq W'\xrightarrow{e}W$
of $f$, with $e$ a weak equivalence and with relative cells of dimension
at least two. Explicitly, attach two-cells along generators of
$\ker f_*$, extending $f$ by their nullhomotopies. Next add spheres to
make each higher homotopy-group map onto and cells one dimension higher
to kill its kernel, proceeding by dimension. At each step use all
needed representatives, and extend their characteristic maps to $W$.
The relative construction fixes the original map on $BH$. The
fundamental group becomes $\pi_1W$ after the first step, and the
successive higher steps preserve the groups already corrected; hence
$e$ is a weak equivalence. The homotopy groups of a $\lambda$-cell CW
complex have at most $\lambda$ elements, since a map from a finite
sphere has image in a finite subcomplex and a finite CW complex has
countable homotopy groups. Thus at most $\lambda$ cells are required.

Apply the successor construction of A.4 to the relative cells of
$(W',BH)$, starting with $BH\xrightarrow{\mathrm{id}}BH$. Its proof
uses only acyclicity of the current map, surjectivity on $\pi_1$, the
size lemma, and the attaching sphere. All these hold at the initial
stage and are preserved at each successor and limit. The construction
embeds the current group at every step and keeps the given classifying
model as a CW subcomplex. At its union it yields an acyclic map
$BG\to W'$ restricting to the inclusion of $BH$. Composing with $e$
proves the assertion and the stated equality on $BH$. The cardinal
calculation is precisely the calculation at successors and limits in
A.4.
\end{proof}

\begin{claim}{A.6}{Proved}{A simple kernel can be imposed in R2}
Let $W$ be connected and $\mathcal U$-small, and let $\mu\in\mathcal U$.
For every infinite cardinal $\lambda\in\mathcal U$ with
$\lambda>\mu$ and with a $\lambda$-cell CW model of $W$, there is an
acyclic map $t:BG\to W$ with
\[
 |G|=|P|=\lambda,\qquad
 P=\ker(\pi_1t:G\twoheadrightarrow\pi_1W)
\]
and $P$ nonabelian simple. Consequently every homomorphism from $G$
to a group of cardinality at most $\mu$ kills $P$.
\end{claim}
\begin{proof}
We first give the relative simple-envelope operation, including its
map to $W$. Suppose $f:BG\to W$ induces an epimorphism
$\tau:G\twoheadrightarrow\pi$, where $\pi=\pi_1W$, and $|G|\leq\lambda$.
Replace $G$ by $G*F_\lambda$, sending each new free letter to $1$ in
$\pi$ and mapping its one-cell to the base point of $W$. At a current
stage $L\twoheadrightarrow\pi$, perform the simple-envelope step of
A.1 for all $g\ne1,h\in\ker(L\to\pi)$, with both new letters
$t_{g,h},s_{g,h}$ sent to $1$ in $\pi$. The words $w$ and $wh$ still
have infinite order by the same free-product normal form, so the HNN
construction embeds $L$. The quotient map to $\pi$ extends because
both sides of every relation have trivial image.

The classifying model of this extension is obtained from the previous
aspherical model by one-cells for the free and stable letters and
two-cells for the HNN relations. Its universal cover is the tree of
contractible spaces from A.3. The existing map to $W$ extends over
these cells: use the constant loop on each new one-cell, and then a
nullhomotopy on each two-cell, since its attaching loop has trivial
image in $\pi_1W$. Choose cellular nullhomotopies and keep all previous
choices fixed.

Iterating countably produces an embedding $G\hookrightarrow S$, a
map $BS\to W$ extending $f$, and an epimorphism $S\to\pi$.
Its kernel $Q$ is simple. Indeed, for $1\ne g,h\in Q$, both elements
belong to the kernel at a common finite stage. The next HNN step gives
$h=w^{-1}s^{-1}ws$ with $w$ in the normal closure of $g$, and all the
conjugating letters in the kernel. Thus $h$ belongs to the normal
closure of $g$ \emph{inside $Q$}. This is the needed stronger assertion;
normal generation merely inside $S$ would not suffice. The kernel
contains the adjoined $F_\lambda$, hence is nonabelian of cardinal
$\lambda$. The whole group has that cardinal as well.

Start now with an acyclic resolution $BG_0\to W$ of size at most
$\lambda$, supplied by A.3 or A.4. Apply the relative simple envelope
and then A.5, alternately:
\[
 G_0\hookrightarrow S_0\hookrightarrow G_1\hookrightarrow S_1
       \hookrightarrow G_2\hookrightarrow\cdots .
\]
Every arrow commutes with the epimorphisms to $\pi$, and their chosen
classifying models and maps to $W$ extend one another literally. The
maps $BG_i\to W$ are acyclic, while
$\ker(S_i\to\pi)$ is nonabelian simple of cardinal $\lambda$.
All groups and models stay within $\lambda$ by A.5.

Take the union of the models. It is aspherical: a sphere and a proposed
nullhomotopy involve finitely many cells and occur in one stage;
all positive homotopy groups above degree one vanish there. Its group
is $G=\operatorname{colim}_iG_i=\operatorname{colim}_iS_i$.
The induced map $t:BG\to W$ is acyclic. For a local system $M$ on $W$,
its cellular homology is the colimit of the homologies of the cofinal
$BG_i$, each mapping isomorphically to $H_*(W;M)$; the transition maps
commute with these isomorphisms. The kernel $P$ is the union of the
cofinal simple kernels $\ker(S_i\to\pi)$. A nontrivial element and
any other element of $P$ lie in a common later simple kernel, where
the first normally generates the second using conjugators in $P$.
Thus $P$ is simple. Its cardinality and that of $G$ are $\lambda$.
The restriction to $P$ of a homomorphism $G\to L$, $|L|\leq\mu$,
has normal kernel in $P$; simplicity makes it trivial or injective,
and cardinality excludes the latter. This proves every assertion.
\end{proof}

\begin{claim}{A.7}{Proved}{The kernel's homology is determined by the target}
For every acyclic map $t:BG\to W$ with $W$ connected and
$P=\ker\pi_1t$, pulling back the universal cover of $W$ gives a
$\pi_1W$-equivariant acyclic map
\[
 BP\longrightarrow\widetilde W.
\]
It induces $\mathbb Z[\pi_1W]$-module isomorphisms
$H_n(P;\mathbb Z)\cong H_n(\widetilde W;\mathbb Z)$ for every $n$.
The action on the left is conjugation through $G/P$.
For fixed $W$, these are exactly the possible graded homology modules
of kernels in R2, including R2 with a nonabelian simple kernel.
\end{claim}
\begin{proof}
The pullback of $\widetilde W\to W$ is the connected cover of $BG$
corresponding to $P$, because $\pi_1t$ is onto. It is aspherical and
has fundamental group $P$. The pullback map has the same homotopy
fibre as $t$, so is acyclic. Its comparison of integral homology is
therefore an isomorphism. Deck transformations commute with that map;
a deck transformation in the covering of $BG$ induces the usual
conjugation action of $G/P$ on $H_*(P)$. This proves the equivariance
and the forced value. Conversely A.6 constructs a simple kernel with
these values for every $W$ and $\mu$. In particular, an independently
prescribed graded module is realized exactly when it is isomorphic to
the displayed graded module. The same pullback map also preserves the
cohomology ring and its cup products.
\end{proof}

\begin{claim}{A.8}{Proved}{Functoriality for injective maps and for all simplicial maps}
Fix an infinite $\lambda\in\mathcal U$. The construction in A.3 can be
chosen as a functor on ordered simplicial complexes and simplicial maps
that are injective and order-preserving on every simplex. Its maps to
geometric realization are natural. Precomposing with barycentric
subdivision gives this assertion for arbitrary simplicial complexes
and all maps injective on every simplex, in particular for injective
simplicial maps.

There is also a functor $\widehat T_\lambda$ on all
$\mathcal U$-small simplicial complexes and all simplicial maps, with a
natural acyclic map $\widehat T_\lambda(K)\to |K|$, aspherical
components, and the no-small-quotient conclusion for groups of
cardinality $<\lambda$. It has at most
$\max(\lambda,\aleph_0,|K|)$ cells. This last functor uses the explicit
resolution below; it does not identify the original simplex cones
with their replacements.
\end{claim}
\begin{proof}
For each $n\geq2$, choose once a group $C_n$ of cardinal $\lambda$
containing $\pi_1T(\partial\Delta^n)$, and choose once the corresponding
map and cylinder $Q_n$. A simplex with its ordered vertices is a
specified copy of $\Delta^n$. Use the same $Q_n$ on every such simplex
and identify its boundary by these ordered face maps. Induction gives
strict face compatibility. A simplexwise injective order-preserving
map carries each copy of $Q_n$ identically to the copy over its image;
the boundary identifications agree, so these maps glue, preserve
composition and identities, and commute with $t$. The proofs of A.3
apply unchanged. Barycentric vertices are nonempty faces, ordered
inside each simplex by inclusion. A simplexwise injective map
preserves their dimensions and strict chains, so gives exactly the
ordered maps just considered.

Here are the replacements needed for a map that collapses vertices.
For a simplicial set $S$, let $\mathcal D_S=\Delta_{\mathrm{inj}}/S$:
its objects are all simplices $(n,s)$, including degenerate ones, and
its morphisms $(m,x)\to(n,s)$ are injective order-preserving maps
$\alpha:[m]\to[n]$ with $x=\alpha^*s$. Every nonidentity arrow
strictly raises $n$. Thus every nondegenerate simplex of
$N\mathcal D_S$ has distinct vertices. Its nondegenerate simplices,
ordered by face inclusion, form a poset $P_S$: a face occurrence is
unique because its vertices determine the subset of the strictly
increasing degree sequence. Put $R(S)=NP_S$, viewed as an ordered
simplicial complex. Its realization is the barycentric subdivision
of $|N\mathcal D_S|$.

A simplicial map $S\to S'$ preserves the dimensions of the objects of
$\mathcal D_S$, since degenerate simplices have not been discarded.
It consequently preserves lengths of nonidentity chains and their
proper face inclusions. The induced map $R(S)\to R(S')$ is therefore
injective on every simplex and order-preserving. These assignments
are functors and do not increase an infinite cardinal bound on the
number of simplices.

We check the comparison to $S$. Include
$\mathcal D_S\to\Delta/S$. Over an object $(n,s)$ of $\Delta/S$, the
comma category has objects all monotone maps $\beta:[m]\to[n]$ and
morphisms injective maps over $[n]$. Its endofunctor appending the
last vertex $n$ has natural transformations from the identity and
from the constant object $([0],n)$. Its nerve is contractible.
Quillen's Theorem A makes this inclusion a nerve weak equivalence.
The last-vertex map $N(\Delta/S)\to S$ is a weak equivalence as well:
the usual simplicial replacement of the diagram
$(n,s)\mapsto\Delta^n$ maps to $S$, and over every $m$-simplex
$x\in S_m$ its indexing category has initial object $(m,x,\mathrm{id})$.
Its projection to $N(\Delta/S)$ is an equivalence because each
$\Delta^n$ is contractible. The two resulting maps to $S$ are
homotopic by their linear homotopy inside the final simplex of each
chain. This proves, with its specified natural last-vertex map,
\[
 |R(S)|\longrightarrow |N\mathcal D_S|\longrightarrow |S|
 \quad\text{is a weak equivalence.}
\]

To pass from an abstract complex $K$ to a simplicial set naturally
for every vertex map, let $U(K)_n$ consist of tuples
$(v_0,\ldots,v_n)$ whose set of entries is a simplex of $K$.
Repetition and arbitrary order are allowed; faces delete entries
and degeneracies repeat them. The affine map
$|U(K)|\to |K|$ sends simplex coordinates to the corresponding
weighted sum of vertices and is strictly natural. For a single
simplex $\sigma$, $U(\sigma)$ is the nerve of the indiscrete category
on its vertices, hence contractible. The affine map is consequently
a weak equivalence on every simplex. Both constructions preserve
unions and intersections of subcomplexes, and the realization
inclusions are cofibrations. Induction over finite subcomplexes,
using the gluing theorem, and then union prove that
$|U(K)|\to|K|$ is a weak equivalence for every $K$.

Define, with the ordered-complex functor from the first paragraph,
\[
 \widehat T_\lambda(K)=T_\lambda\bigl(R(U(K))\bigr).
\]
The displayed natural comparisons give an actual natural map to
$|K|$. It is acyclic by A.3 followed by weak equivalences, and its
components are aspherical. Its fundamental-group kernels are the
kernels in A.3, transported by the fundamental-group isomorphisms
of the comparisons, so the no-small-quotient property holds. The
constructions use finite strings and countable unions only, giving
the asserted cell bound. Inclusions of complexes stay inclusions
through this construction and therefore still inject on fundamental
groups by A.3.
\end{proof}

\begin{claim}{A.9}{Proved}{The limits of simultaneous functorial strengthening}
The original simple-cone construction, with
$T(\Delta^n)\simeq BC_n$ and subcomplex group injections, cannot extend
to a functor on all ordered simplicial maps while keeping those
values and maps. Moreover, no acyclic-resolution functor on all
connected spaces and maps can have a nonabelian simple kernel on
every simply connected space.
\end{claim}
\begin{proof}
For the first assertion take a face inclusion
$i:\Delta^2\hookrightarrow\Delta^3$ and its simplicial retraction
$r:\Delta^3\to\Delta^2$, obtained by identifying the remaining
vertex with an endpoint. Functoriality would give a retraction
$C_3\to C_2$ of the already injective face map. The retraction is
nontrivial and $C_3$ is simple, so it is injective as well as onto.
Consequently the face subgroup is all of $C_3$. But
$\pi_1T(\partial\Delta^3)$ embeds into $C_3$ and contains that face
subgroup, so it too would be isomorphic to $C_3$. The specified
homology comparisons of A.3 give
\[
 H_2(T(\partial\Delta^3);\mathbb Z)=\mathbb Z,
 \qquad H_2(BC_3;\mathbb Z)=0,
\]
a contradiction. This proves why A.8 supplies a changed model for
maps collapsing vertices.

For the second assertion use the inclusion
$i:S^2\to S^2\vee S^2$ of the first summand and the retraction $r$
collapsing the other summand. Both targets are simply connected,
so the hypothetical resolving groups $G_1,G_2$ themselves are
nonabelian simple. Their homomorphisms would have
$G(r)G(i)=\mathrm{id}$. The map $G(r)$ is nontrivial, since its
map on $H_2$ must be the projection
$\mathbb Z^2\to\mathbb Z$ under the natural acyclic comparisons.
Simplicity makes $G(r)$ injective, and the retraction makes it onto.
It would induce an isomorphism on $H_2$, contradicting that specified
projection. The same argument applies to a functor into based
aspherical spaces by taking fundamental groups; in the unpointed
version the equations are up to conjugacy, which yields the same
retraction and homology contradiction. The finite cellular
homology computations are independently checked in
\texttt{code/a\_finite\_checks.py}.
\end{proof}

\begin{claim}{A.10}{Proved}{Dimension preservation for all finite target groups}
Let $W$ be a connected CW complex of finite dimension $n$. There is
an acyclic map $t:BG\to W$ with a classifying model of dimension at
most $n$, such that every homomorphism $G\to F$ to a finite group
kills $\ker\pi_1t$. Thus the same resolution satisfies R2 for every
finite cardinal $\mu$. A finite, respectively countable, $W$ admits
a finite, respectively countable, resolving model. More generally,
an infinite bound on the cells of $W$ also bounds the cells of this
model and the cardinality of $G$.
\end{claim}
\begin{proof}
We record the additional group-theoretic verification needed to use
the dimension-preserving basic realization of
\cite[Sections 9--10]{A_BDH1980}. Use the Higman group
\[
 H=\langle a,b,c,d\mid a^{-1}ba=b^2,
 b^{-1}cb=c^2,\ c^{-1}dc=d^2,\ d^{-1}ad=a^2\rangle.
\]
Its presentation complex is a finite aspherical two-complex. Indeed,
the three-generator group on $a,b,c$ is the amalgam of the two
Baumslag--Solitar groups on $(a,b)$ and $(b,c)$ over $\langle b\rangle$;
their standard mapping-torus models are aspherical. The subgroup
$\langle a,c\rangle$ is free of rank two by amalgam normal form:
nonzero powers of $a$ and $c$ do not lie in the edge group.
The analogous group on $c,d,a$ has the same free subgroup, and their
amalgam over it is $H$. The tree-of-spaces argument proves the asserted
asphericity and embeds these groups. The cellular map from its four
two-cells to its four one-cells is a signed permutation matrix, because
the abelianized relators have respective exponent sums $-b,-c,-d,-a$.
It is an isomorphism, so the complex is integrally acyclic.

Every homomorphism from $H$ to a finite group is trivial. Otherwise
choose the smallest prime $p$ dividing the order of an image of one
of the four generators, say $x$. The preceding generator $y$
conjugates $x$ to $x^2$. Thus the order of $x$ is odd, and the
multiplicative order of $2$ modulo $p$ divides the order of $y$.
This multiplicative order is greater than one and divides $p-1$;
one of its prime divisors is smaller than $p$, contradicting the
choice of $p$. This proves the required property of $H$ without
assuming simplicity of that group.

Here are the relevant details of the BDH construction, in order to
check the quotient property not stated in its theorem. Write $L(K)$
and $\widehat L(K)$ for its unpointed and pointed groupoids. In
dimension zero, $L$ is discrete, while $\widehat L$ adds one common
base object and its unique arrows to the vertices. For an
$m$-simplex $\sigma$, the construction sets
$L(\sigma)=\beta\widehat L(\partial\sigma)$, where $\beta$ removes
the base object, and defines $\widehat L(\sigma)$ by the pushout
of groupoids
\[
 \varphi\widehat L(\partial\sigma)
 \ \longleftarrow\ L(\partial\sigma)
 \ \longrightarrow\ H\times L(\partial\sigma).
\]
Here $\varphi$ forgets the distinguished object. Both arrows are
injective, and the new map from $L(\sigma)$ is obtained by conjugating
an arrow $u:x\to y$ by $(h,y)^{-1}$ and $(h,x)$ for a fixed
$1\ne h\in H$. The reduced-word calculation of BDH Lemma 10.1
shows that the resulting boundary maps and all subcomplex maps are
injective. Their Theorem 9.1 identifies each groupoid pushout with
a homotopy pushout of classifying spaces.

In the pointed groupoids, every vertex group is normally generated
by the inserted copies of $H$. To see this directly, kill those
copies in the displayed pushout. The factor $H\times
L(\partial\sigma)$ then becomes $L(\partial\sigma)$, and the
pushout reduces to $\varphi\widehat L(\partial\sigma)$. Its quotient
is the indiscrete groupoid on its objects by induction. Gluing
pointed simplices adds no loops to this quotient: every intersection,
including the intersection of disjoint simplices, contains their
common base object and has an indiscrete quotient. The union quotient
is therefore again indiscrete. This proves normal generation, for
finite complexes and then for arbitrary unions.

Consequently every homomorphism from a vertex group of
$L(\sigma)=\beta\widehat L(\partial\sigma)$ to a finite group is
trivial. The kernel of the specified comparison
$\pi_1BL(K)\to\pi_1|K|$ is exactly the normal closure of these
simplex vertex groups: after they are killed, each simplex groupoid
is indiscrete, and their colimit is the fundamental groupoid of
$K$ by van Kampen. Every finite-group homomorphism kills this
kernel. BDH Lemma 10.2 and Proposition 8.13 prove the comparison
with all local coefficients (a simplex carries a constant pulled-back
system). Thus this is the required acyclic map.

For clarity, the finite-dimensional assertion does not follow just
by taking an unspecified union of low-dimensional models. BDH
Theorem 10.4 gives a finite model of dimension at most $m$ for the
acyclic groupoid $L(\Delta^m)$, for each finite standard simplex.
Construct compatible models over $K$ by dimension: a graph in
dimensions zero and one, and over every $m$-simplex attach the
mapping cylinder from the already chosen boundary model to that
$m$-dimensional classifying model, using the specified injective
group homomorphism. The boundary model has dimension at most
$m-1$, so its cylinder has dimension at most $m$. The same
normal-form, local-coefficient, and subcomplex arguments as in A.3
identify these models with the prescribed groupoids. Their union
is an actual $n$-dimensional aspherical CW complex. There are
finitely many new cells for each simplex, yielding the stated
finite and cardinal bounds. Finally replace $W$ by the
dimension-preserving simplicial model from A.4 and compose with
the homotopy equivalence to $W$.
\end{proof}

\begin{claim}{A.11}{Proved}{The exact cardinal restriction for countable models}
For any resolution $t:BG\to W$ with nontrivial kernel $P$, the
no-small-quotient assertion of R2 at $\mu$ implies $\mu<|G|$.
For countable finite-dimensional nonaspherical $W$, a countable
resolving model of dimension at most $\dim W$ satisfies R2 at
$\mu$ precisely for finite $\mu$, in the sense of existence of
such a model. Dimension alone imposes no bound on $\mu$ valid for
all $W$: the identity map of any connected graph satisfies R2 for
every $\mu$ and has the same dimension as its target.
\end{claim}
\begin{proof}
If $|G|\leq\mu$, the homomorphism $\mathrm{id}_G:G\to G$ is one
of the tested homomorphisms and does not kill $P$. This proves the
necessary cardinal inequality. If an acyclic map $BG\to W$ has
$P=1$, it is an isomorphism on fundamental groups and on homology
with the group-ring local system; the universal-cover Whitehead
argument of A.2 then makes it a weak equivalence. Thus a
nonaspherical $W$ forces $P\ne1$. A countable CW model has a
countable fundamental group, so the inequality rules out every
infinite $\mu$. A.10 provides such a countable model working for
every finite $\mu$. Finally, a connected graph is aspherical and
its identity has trivial kernel, proving the last assertion by the
specified map. These conclusions concern cardinality of the
source; no conversion from dimension to cardinality is used.
\end{proof}

